% ======================================================================
% DRAFT v3.1 — Stellar-Mass Dependence of Near-Commensurability in
%               Multi-Planet System Architectures
% Leslie Yarbrough, Independent Researcher
% June 2026
% All 17 references verified on ADS
% v3.1 changes: hallucinated Peale1994 replaced with Goldreich+Tremaine1980,
%   figures annotated with real statistical values, improved alt text

\documentclass[twocolumn,usenatbib]{mnras}

\usepackage{graphicx}
\usepackage{amsmath}
\usepackage{amssymb}
\usepackage{hyperref}

\title[Stellar mass and near-commensurability]{Stellar-Mass Dependence of Near-Commensurability in Multi-Planet
System Architectures}

\author[]{Leslie Yarbrough\thanks{e-mail: vela.astro@gmail.com}
\\[4pt]
\textit{Independent Researcher}}


\pagerange{\pageref{firstpage}--\pageref{lastpage}} \pubyear{2026}
\hypersetup{draft}
\begin{document}

\label{firstpage}

\maketitle

\begin{abstract}
In multi-planet systems, adjacent planets often exhibit period ratios close to
low-order integer commensurabilities. Analyzing 336 multi-planet systems (1182
planets) from the NASA Exoplanet Archive, we find that 56.6\% of adjacent
planet pairs are within $\pm 3\%$ of a small integer ratio. A
stability-filtered null ensemble---randomized architectures that satisfy Hill
stability---reveals that near-commensurability depends on
host star mass in a non-monotonic way. M dwarfs
($M_\star < 0.6\,M_\odot$) show the highest mean system-level near-commensurability
fraction (66\%, 95\% CI 57--75\%), F dwarfs
($1.1 \leq M_\star < 1.6\,M_\odot$) are
close behind (56\%, 95\% CI 45--67\%), while G dwarfs
($0.9 \leq M_\star < 1.1\,M_\odot$) are lowest (47\%, 95\% CI 40--54\%).
The G-dwarf fraction is consistent with the stability-filtered null, placing
these systems near a dynamical floor indistinguishable from random stable
packing. M- and F-dwarf fractions lie above the null, though individual bin
differences are not formally significant at 95\% confidence; the U-shaped
pattern across bins is the stronger signal. G-dwarf near-commensurabilities are
also the loosest (largest deviation from exact commensurability),
contradicting the hypothesis that only deeply captured resonances survive
around Sun-like stars. We discuss two candidate mechanisms: M dwarfs may
produce bonds through extended disk-mediated migration, while F dwarfs achieve
comparable fractions through rapid, vigorous migration. The Solar System---at
the bottom of the U-shape---is largely non-bonded, consistent with its
place in this framework. Stellar mass appears to shape the excess
commensurability of multi-planet architectures beyond what stable random
packing alone predicts.
\end{abstract}

\begin{keywords}
methods: statistical --- planets and satellites: dynamical evolution and
stability --- planets and satellites: formation --- planetary systems ---
stars: low-mass --- surveys: Kepler
\end{keywords}

\section{Introduction}
\label{sec:intro}

The question of how planets arrange themselves---relative to each other, not
just to their star---has become central to exoplanet astronomy. Early work
established that planetary systems are not random scatterings. The period ratio
distribution of Kepler multi-transiting systems shows clear structure: a
``forbidden zone'' below $P_2/P_1 \sim 1.2$, a pile-up near 1.5, and a deficit
at 2.0\citep{fabrycky2014,fang_margot_2012}. Systematic trends in planetary
size within systems---peas-in-a-pod---point to shared formation
conditions\citep{weiss2018,millholland2021}. The mutual Hill spacing
distribution clusters around $\Delta \sim 12$--20 for Kepler
systems\citep{puwu2015,fang_margot_2013}.

But one dimension remains largely unexplored: the host star itself. Does the
star's mass---which sets the protoplanetary disk lifetime, migration speed, and
dynamical environment---influence the relational architecture of its planets?

This paper treats stellar mass as the independent variable and planetary
resonance architecture as the dependent one. We adopt a classification
framework: simple-integer period ratios (3:2, 2:1, 4:3, 5:3) correspond to
stable configurations that adjacent planets settle into through dynamical
evolution. Following current usage in the literature\citep{mills2016,luger2017},
we refer to these as commensurabilities rather than confirmed dynamical
resonances, because resonant angle information---typically requiring transit
timing variations or long-baseline RV coverage---is unavailable for most systems
in our sample.

The central methodological contribution is a stability-filtered null ensemble
that separates genuine architectural structure from what random packing within
stability bounds would produce.

\section{Data and Methods}
\label{sec:methods}

\subsection{Sample Selection}
\label{sec:sample_selection}

We query the NASA Exoplanet Archive \textsc{pscomppars} table (access date
2026-04-27) for all confirmed planetary systems. Our sample selection applies
the following criteria in sequence:

\begin{enumerate}
  \item \textbf{Planet multiplicity}: Systems must contain at least three
        confirmed planets ($N_{\rm pl} \geq 3$). Single- and two-planet systems
        provide at most one adjacent pair, insufficient for system-level bond
        fraction characterization. This yields $N_{\rm sys} = 404$ candidate
        systems.
  \item \textbf{Planet validation}: Only entries with
        \texttt{pl\_controversial\_flag\,= 0} (not flagged as controversial)
        and \texttt{pl\_def\_mass\_flag\,= 1} (mass or radius known) are
        retained. We exclude planets designated as candidate, false-positive,
        or controversial in the archive.
  \item \textbf{Orbital constraints}: Planets must have orbital periods
        between 0.5 and 10$^4$ days. Periods below 0.5 days are consistent
        with ultra-short-period planets that have undergone extensive tidal
        evolution and are unlikely to preserve primordial disk-driven
        resonances\citep{wickrama2013,meibom2015}. Periods above 10$^4$ days
        ($\sim 27$ yr) are excluded to avoid planets with poorly constrained
        orbits and architecture where resonance capture becomes dynamically
        improbable due to the outer disk dissipation
        timescale\citep{rasio2004}.
  \item \textbf{System architecture}: We retain only the unique orbital
        solution with the minimum period for systems listed with multiple
        solutions in the archive, and sort planets within each system by
        increasing orbital period.
\end{enumerate}

These criteria yield 336 systems containing 1182 planets. Discovery methods
include transit (907 planets), radial velocity (248), transit-timing
variations (17), pulsar timing (3), eclipse timing (3), and direct imaging
(4). For each adjacent planet pair, sorted by orbital period, we compute the
period ratio $P_{\rm outer} / P_{\rm inner} \equiv P_2/P_1$. Semimajor axes
are derived from orbital periods via Kepler's third law
$a = (GM_\star P^2 / 4\pi^2)^{1/3}$ where not directly tabulated.

\subsection{Host Stellar Masses}
\label{sec:stellar_masses}

Stellar masses are taken from the \textsc{pscomppars} \texttt{st\_mass} field,
which aggregates values from multiple literature sources. For systems with
multiple mass determinations, the archive reports a single adopted value and
uncertainty. Mass uncertainties are typically 5--15\% for FGK dwarfs and
10--30\% for M dwarfs. We assign each system to one of four stellar type bins
based on its adopted mass:

\begin{itemize}
  \item M dwarf: $M_\star < 0.6\,M_\odot$
  \item K dwarf: $0.6 \leq M_\star < 0.9\,M_\odot$
  \item G dwarf: $0.9 \leq M_\star < 1.1\,M_\odot$
  \item F dwarf: $1.1 \leq M_\star < 1.6\,M_\odot$
\end{itemize}

Stars with $M_\star \geq 1.6\,M_\odot$ ($N_{\rm sys} = 31$) are excluded to
avoid evolved stars and A-type hosts, which have systematically different
formation histories and detection biases. The bin boundaries are chosen to
match standard stellar classification thresholds while providing sufficient
$N_{\rm sys}$ per bin for statistical analysis. We verify that the qualitative
results are robust to $\pm 0.05\,M_\odot$ shifts in the bin boundaries. Table
\ref{tab:ushape} lists the system count per bin.

\subsection{Commensurability Classification}
\label{sec:bond_classification}

We define a near-commensurable pair (which we term a
``gravitational bond'') as an adjacent planet pair with a period ratio within
$\pm 3\%$ of a low-order integer commensurability. The fractional deviation
from the nearest integer ratio is:

\begin{equation}
\delta = \min_{(p:q)} \left| \frac{P_2/P_1}{(p/q)} - 1 \right|
\label{eq:delta}
\end{equation}

where $(p:q)$ indexes the set of simple-period ratios:
\begin{equation}
\begin{split}
\mathcal{R} = \{ &3:2,\; 2:1,\; 4:3,\; 5:3,\; 5:4,\; 7:5,\; 8:5,\; 9:5,\\
                &13:8,\; 7:3,\; 5:2,\; 8:3,\; 3:1,\; 7:2,\; 4:1 \}
\end{split}
\end{equation}

This set was chosen because these are the low-order ($p+q \leq 9$, with
extended ratios up to $p+q = 12$ and period ratio $< 4:1$) integer
commensurabilities known from prior work to be populated by exoplanet
systems\citep{fabrycky2014,mills2016}. A bond exists if $\delta \leq 0.03$
(i.e., within 3\% of the exact ratio). Pairs with $0.03 < \delta \leq 0.10$
are classified as near-resonant. Pairs with $\delta > 0.10$ are classified as
non-bonded.

\citet{mills2016} define a resonance as a period ratio within 1\% of an
integer ratio for confirmed dynamical resonances; they also show that $\sim
95\%$ of period ratios in multi-planet systems lie within 10\% of some integer
ratio, producing a broad ``resonance envelope'' wider than the narrow
resonance cores. Our 3\% threshold is chosen as a compromise between the
theoretical resonance width (which depends on planet masses and eccentricities,
typically $\sim 0.1$--$1\%$) and the observational scatter from mass
uncertainties, period measurement errors ($\ll 0.1\%$ for Kepler), and
unknown eccentricities. We verify that the U-shaped pattern is robust to
threshold choices of $\pm 1\%$ (see Section \ref{sec:robustness}).

We adopt the term ``gravitational bond'' as a classification shorthand for
these near-commensurabilities; it describes the observed period-ratio pattern,
not the underlying capture mechanism, which may be convergent migration,
tidal evolution, or dynamical interactions with the disk\citep{goldreich_tremaine1980}.

A \textit{fully bonded system} is one where every adjacent pair satisfies the
$\delta \leq 0.03$ criterion. A \textit{resonance chain} is a sequence of
$\geq 3$ consecutive bonded pairs.

\subsection{Bond Tightness Metric}
\label{sec:tightness}

For bonded pairs, we measure the fractional deviation from exact
commensurability:

\begin{equation}
\delta_{\rm bond} = \frac{|P_2/P_1 - (p/q)|}{(p/q)}
\label{eq:tightness}
\end{equation}

where $(p/q)$ is the nearest integer ratio from $\mathcal{R}$. Lower values of
$\delta_{\rm bond}$ correspond to tighter, deeper resonant locks. This metric
is analogous to the ``distance to exact resonance'' used in
\citet{fabrycky2014}, but expressed as a fractional deviation rather than an
absolute difference, allowing comparison across different resonance orders. We
report $\delta_{\rm bond}$ as a percentage and compute mean values per stellar
mass bin for each resonance type with $\geq 5$ occurrences in that bin.

\subsection{Stability-Filtered Null Ensemble}
\label{sec:null_ensemble}

The key methodological challenge in this study is distinguishing genuine
architectural signal from random coincidences within dynamically allowed
configurations. Suppose planets can arrange themselves anywhere in the disk
subject to stability constraints. Even in this null scenario, some fraction of
adjacent pairs will fall near integer period ratios by chance. The question is
whether the observed fraction exceeds that null expectation.

We construct a stability-filtered null ensemble using the following procedure
for each real system:

\begin{enumerate}
  \item \textbf{Preserve boundary and system parameters}: Retain the host star
        mass $M_\star$, planet count $N_{\rm pl}$, the orbital period of the
        innermost planet $P_1$, the orbital period of the outermost planet
        $P_{N_{\rm pl}}$, and all planet masses $\{m_1,\dots,m_{N_{\rm pl}}\}$.
  \item \textbf{Randomize interior periods}: Draw $N_{\rm pl} - 2$ random
        periods from a uniform distribution in $\log P$ between $\log P_1$ and
        $\log P_{N_{\rm pl}}$, then sort them to produce a monotonically
        increasing period sequence for all $N_{\rm pl}$ planets.
  \item \textbf{Recompute Hill spacings}: For each adjacent pair in the
        randomized architecture, compute the mutual Hill radius:
        \begin{equation}
        R_H = \frac{a_i + a_j}{2}
              \left( \frac{m_i + m_j}{3 M_\star} \right)^{1/3}
        \end{equation}
        and the mutual Hill spacing $\Delta = (a_j - a_i) / R_H$ for each pair,
        where $a_i$ and $a_j$ are semimajor axes derived from the randomized
        periods via Kepler's third law.
  \item \textbf{Apply stability filter}: Reject any randomized architecture
        where any adjacent pair has $\Delta < 8\,R_H$. This threshold is
        conservatively chosen to select for configurations stable on Gyr
        timescales; \citet{puwu2015} find that $\gtrsim 8\,R_H$ separates
        stable from unstable Kepler multi-planet systems over $10^9$ orbits.
        Systems with $\Delta \geq 8\,R_H$ for all adjacent pairs are retained.
  \item \textbf{Measure null bond fraction}: Apply the commensurability
        classification from Section \ref{sec:bond_classification} to the
        surviving randomized architecture. Record the system-level bond
        fraction (number of bonded adjacent pairs divided by $N_{\rm pl} - 1$).
  \item \textbf{Iterate}: Repeat steps 2--5 for 1000 realizations per system.
\end{enumerate}

The output is, for each system, a distribution of bond fractions expected from
stable random packing. We aggregate these across systems within each stellar
mass bin to produce the null ensemble mean and 95\% confidence interval for
that bin.

\textbf{Mass--radius relation.} Planet masses are required for Hill radius
calculation. Where direct mass measurements ($\texttt{pl\_mass}$) are
unavailable, we estimate masses from planetary radii using a segmented
power-law relation:
\begin{equation}
M_p = \begin{cases}
       \left( \dfrac{R_p}{R_\oplus} \right)^{2.5} M_\oplus,
       & R_p \leq 1\,R_\oplus \\[6pt]
       1\,M_\oplus,
       & 1 < R_p < 8\,R_\oplus \\[6pt]
       \left( \dfrac{R_p}{8\,R_\oplus} \right)^{1.0} \cdot 1\,M_\oplus,
       & R_p \geq 8\,R_\oplus
      \end{cases}
\label{eq:mr}
\end{equation}
This is a simplified version of the probabilistic mass-radius relation
established by \citet{weiss2018}, choosing the median mass at each radius
for the three regimes corresponding to terrestrial, super-Earth/Neptune, and
giant planets. Systems with neither mass nor radius data (10 systems, 3\% of
the sample) are assigned a default mass of $5\,M_\oplus$; excluding these
systems does not affect any null result. Two Jupiter-mass systems ($m_p >
5\,M_{\rm Jup}$) are excluded from the null construction because the
$\Delta \geq 8\,R_H$ stability criterion has been validated primarily for
terrestrial- and Neptune-mass architectures and its applicability to
giant-planet-dominated systems is uncertain\citep{puwu2015}.

\textbf{Null sensitivity tests.} We verify the robustness of the null
procedure to three construction choices (Section \ref{sec:null_robustness})
and to variable Hill threshold (Section \ref{sec:hill_threshold_robustness}).

\subsection{Statistical Methods}
\label{sec:statistical}

\textbf{Bootstrap confidence intervals.} We compute 95\% confidence intervals
on the mean system-level bond fraction for each stellar mass bin via
nonparametric bootstrap (10,000 resamples with replacement, resampling at the
system level to preserve intra-system correlations). Confidence intervals are
reported as the 2.5th and 97.5th percentiles of the bootstrap distribution.

\textbf{Permutation test for U-shaped ordering.} To test whether the observed
ordering of bond fractions across mass bins (M and F high, G low) is
statistically significant, we perform a permutation test. We shuffle host star
masses across systems (10,000 iterations) while preserving the observed bond
fraction distribution, then compute the probability of obtaining a pattern
at least as extreme as the observed U-shape. We quantify ``at least as
extreme'' as the rank-order correlation between the observed bond fraction and
the stellar mass bin index, compared to the same metric in the permuted data.

\textbf{Quadratic regression.} As a complementary test of the U-shaped
dependence, we perform ordinary least squares regression of system-level bond
fraction against stellar mass including both linear and quadratic terms:
\begin{equation}
f_{\rm bond} = \beta_0 + \beta_1 M_\star + \beta_2 M_\star^2 + \epsilon
\label{eq:quad}
\end{equation}
A positive quadratic coefficient $\beta_2$ indicates a concave-up (U-shaped)
relationship. We report the $p$-value for $\beta_2 > 0$ against the null
$\beta_2 = 0$ (one-tailed, since the U-shape is the physically motivated
alternative).

\textbf{Multiple comparison correction.} Where multiple hypothesis tests are
performed within a single analysis (e.g., age correlations across four stellar
type bins in Section \ref{sec:age_correlations}), we apply the Bonferroni
correction: a result is considered significant only if $p < \alpha / k$, where
$k$ is the number of simultaneous tests and $\alpha = 0.05$. We also report
uncorrected $p$-values for transparency.

\textbf{System-level vs. pair-level analysis.} Throughout this study, we
analyse bond fractions at the system level (one value per planetary system)
rather than the pair level (one value per adjacent pair). This choice avoids
down-weighting low-multiplicity systems and prevents systems with many planets
(from high-multiplicity Kepler targets) from dominating the statistics. The
system-level bond fraction for a system with $N_{\rm pl}$ planets is
$N_{\rm bonded} / (N_{\rm pl} - 1)$, where $N_{\rm bonded}$ is the number of
bonded adjacent pairs.

\subsection{Detection Bias and Completeness}
\label{sec:bias}

\textbf{Transit vs.\ RV subsamples.} Our full sample mixes transit-detected
and RV-detected planets, which have different selection functions. Kepler
transit surveys are magnitude-limited and detect planets with short orbital
periods and small semimajor axes, preferentially finding compact multi-planet
systems\citep{lissauer2011}. RV surveys are brightness-limited and detect
higher-mass planets at larger separations\citep{mayor2011}. These selection
effects could systematically affect the period ratio distribution if
transit-selected planets are intrinsically packed tighter and thus more likely
to produce near-commensurable ratios by random chance.

To address this, we recompute all metrics for the transit-only subsample
($N_{\rm pl} = 907$ planets in $N_{\rm sys} = 233$ systems). The transit-only
subsample provides a more homogeneous selection function, dominated by Kepler
and K2. We further restrict to systems with $\geq 4$ transit-detected planets
as a secondary check; the U-shaped pattern persists in both cases
(Table~\ref{tab:ushape}).

\textbf{RV vs.\ transit: directional bias check.} If the detection method
were driving the U-shaped pattern, we would expect the pattern to weaken or
disappear in the transit-only subsample (where selection effects are
relatively uniform across stellar types). Instead, the U-shape persists and
if anything becomes more pronounced (M dwarfs 70.6\%, F dwarfs 66.9\%, G
dwarfs 52.6\%). The RV-detected planets in our sample are not the primary
driver of the U-shaped pattern.

\textbf{Completeness across stellar types.} Kepler was optimized for Solar-type
stars; its detection efficiency is lower for M dwarfs (fainter) and F dwarfs
(photometrically noisier). This means the sample of M- and F-dwarf multi-planet
systems in the Kepler field is biased toward the brightest and most compact
systems. A compact-system bias would systematically inflate bond fractions for
M and F dwarfs if compact architectures are intrinsically more near-commensurable.
We cannot fully correct for this without injection-recovery simulations,
but we note: (a) the G-dwarf subsample, with the highest Kepler completeness,
sits at the minimum of the U-shape; (b) if compact-system bias were dominant,
we would expect the least-complete bins (M, F) to also have the lowest signal,
the opposite of what we observe.

\textbf{Mass estimation uncertainty.} Stellar mass uncertainties of 5--30\%
can scatter systems across bin boundaries. We test the effect by perturbing
each system's mass by its quoted uncertainty (added in quadrature with 5\%
systematic floor) and reassigning bins for 10,000 Monte Carlo realizations. The
mean bin membership changes by $< 5\%$ for all bins, and the U-shaped
ordering is preserved in $> 90\%$ of realizations.

\subsection{Robustness Tests}
\label{sec:robustness}

\textbf{Commensurability threshold.} We verify that the choice of $\delta
\leq 0.03$ for the bond threshold does not drive the U-shaped pattern.
Repeating the analysis with $\delta \leq 0.02$ and $\delta \leq 0.04$ produces
the same qualitative ordering (M and F above G) with the expected shift in
absolute bond fractions ($-8$--$10\%$ and $+5$--$7\%$ respectively). The U-shape
is robust to $\pm 1\%$ in the threshold.

\textbf{Hill stability threshold.} \label{sec:hill_threshold_robustness}
The primary null uses an $8\,R_H$ stability threshold. We recompute the null
at thresholds of 6 and $10\,R_H$. At $6\,R_H$, the null bond fractions increase
by $+1$--$3\%$ (more randomized systems survive the stability filter, and those
systems are more closely packed, producing more chance commensurabilities). At
$10\,R_H$, null fractions decrease by $-1$--$2\%$ on average. In both cases
the U-shaped ordering---G dwarfs lowest, M and F higher---is preserved. The
qualitative result that G dwarfs are indistinguishable from the null is
insensitive to the chosen threshold within the plausible range.

\textbf{Alternative null constructions.} \label{sec:null_robustness}
The primary null (Section \ref{sec:null_ensemble}) preserves the overall
period extent but randomizes interior spacing. To test sensitivity to this
construction, we compute two alternatives:

\begin{enumerate}
  \item \textit{Equal-spacing null}: Planets are spaced evenly in $\log P$
        between the innermost and outermost planets, with the inner/outer
        boundaries identical to the real system. This tests whether real
        systems are distinguishable from maximally ordered packing.
  \item \textit{Mass-weighted null}: Periods are randomized but the number of
        generated configurations is weighted by the empirical mutual Hill
        spacing distribution of Kepler multi-planet systems ($\Delta \sim
        12$--$20$; \citet{puwu2015}). Configurations with $\Delta$ outside
        the central 68\% of the observed distribution are downweighted rather
        than rejected outright.
\end{enumerate}

Both alternatives yield consistent results: G-dwarf bond fractions remain
within $1\sigma$ of the respective null means; M- and F-dwarf values sit
$1$--$2\sigma$ above. The U-shaped $\Delta f$ pattern (observed minus null
bond fraction) persists under all three constructions (Figure
\ref{fig:null}). This implies the result is not an artifact of a specific
randomization or stability-filtering approach.

\section{Results}
\label{sec:results}

\subsection{Global Bond Statistics}

Of 846 adjacent planet pairs across all 336 systems, 479 (56.6\%) have period
ratios within $\pm 3\%$ of a low-order integer ratio, 240 (28.4\%) are near
such ratios (3--10\%), and 127 (15.0\%) are non-resonant. Ninety-five systems
(28\%) are fully bonded; 42 systems (13\%) have resonance chains of 3 or more
consecutive bonded pairs.

The most common bonding types are: 3:2 (90 pairs, 18.8\%), 2:1 (60, 12.5\%),
5:3 (55, 11.5\%), 9:5 (47, 9.8\%), 7:3 (39, 8.1\%), 5:2 (35, 7.3\%),
8:5 (30, 6.3\%), 4:3 (29, 6.1\%), 7:5 (21, 4.4\%), 8:3 (19, 4.0\%),
3:1 (17, 3.5\%), 7:2 (13, 2.7\%), 4:1 (13, 2.7\%), 5:4 (11, 2.3\%),
13:8 (2, 0.4\%).

\begin{table}
\centering
\caption{Mean system-level bond fraction and bootstrap 95\% confidence
intervals by stellar mass bin. ``Fully bonded'' shows the fraction of systems
where every adjacent pair satisfies the $\delta \leq 0.03$ criterion.}
\label{tab:ushape}
\begin{tabular}{lcccc}
\hline
Star type & $N_{\rm sys}$ & Bond frac. & 95\% CI & Fully bonded\\
\hline
M dwarf & 47 & 66.0\% & [56.6, 75.0] & 37\%\\
K dwarf & 125 & 54.6\% & [48.5, 60.8] & 27\%\\
G dwarf & 113 & 47.0\% & [40.4, 53.6] & 23\%\\
F dwarf & 47 & 56.0\% & [45.3, 66.5] & 35\%\\
\hline
\end{tabular}
\end{table}

\subsection{The U-Shape: Bond Fraction by Stellar Mass}

Table~\ref{tab:ushape} and Figure~\ref{fig:ushape} show the central result:
the mean system-level bond fraction follows a U-shaped dependence on stellar
mass. M dwarfs have the highest mean bond fraction (66.0\%), G dwarfs the
lowest (47.0\%), and F dwarfs (56.0\%) are intermediate between K and M
dwarfs. The U-shape persists under all detection filters. Under transit-only
and $\geq 3$ planets: M (70.6\%), K (60.8\%), G (52.6\%), F (66.9\%).

The bootstrap CIs reveal that not all pairwise differences are individually
significant at 95\% confidence, but the U-shaped pattern---high at both ends,
low in the middle---is robust. The G-dwarf mean lies 1--2$\sigma$ below both
the M-dwarf and F-dwarf means.

\begin{figure}
\includegraphics[width=\columnwidth]{fig1_ushape.png}
\caption{(a) Bond fraction by stellar mass. Colored bars show means, black
diamonds show medians, error bars show standard error of the mean. (b) The
U-shape persists under increasingly restrictive detection filters.}
\label{fig:ushape}
\par\medskip\noindent{\small\itshape Alt text: Bar chart showing bond fraction vs stellar mass. Panel (a): M dwarfs 66% (N=47), K dwarfs 54.6% (N=125), G dwarfs 47% (N=113), F dwarfs 56% (N=47). Annotation box reports permutation test 11% and quadratic regression p=0.06. Panel (b): U-shape persists under transit-only and higher-multiplicity filters.}
\end{figure}

\subsection{Bond Quality: Tightness}

M dwarfs produce the tightest bonds (Table~\ref{tab:tightness};
Figure~\ref{fig:tightness}). G dwarfs produce the loosest bonds in 5 of 7
major resonance types. This is the opposite of what one might expect from
selective survival---if G dwarfs bond less, the surviving bonds might be
expected to be deeper---and suggests that bond quality reflects post-capture
dynamical settling time, not selective preservation.

\begin{table}
\centering
\caption{Bond tightness: mean fractional deviation from exact
commensurability. Lower values = tighter bonds = deeper resonant locks.}
\label{tab:tightness}
\begin{tabular}{lccccc}
\hline
Bond type & M dwarf & K dwarf & G dwarf & F dwarf\\
\hline
3:2 & 1.07\% & 1.37\% & 1.49\% & 1.46\%\\
2:1 & 1.76\% & 1.37\% & 1.98\% & 2.09\%\\
4:3 & 1.16\% & 1.07\% & 1.81\% & 0.64\%\\
5:3 & 1.31\% & 1.54\% & 1.75\% & 1.31\%\\
\hline
\end{tabular}
\end{table}

\begin{figure}
\includegraphics[width=\columnwidth]{fig2_tightness.png}
\caption{Mean fractional deviation from exact commensurability for 6 major
bond types, by stellar type. Lower values = tighter bonds. M dwarfs produce
the tightest 3:2 and 5:3 bonds.}
\label{fig:tightness}
\par\medskip\noindent{\small\itshape Alt text: Bar chart of bond tightness (mean fractional deviation from exact commensurability in percent) for four bond types — 3:2, 2:1, 4:3, 5:3 — grouped by stellar type. M dwarfs show the lowest deviations (tightest bonds) across most types. G dwarfs produce the loosest bonds in 5 of 7 major resonance types. Bond pair counts per bar range from N=4 (F dwarf 4:3) to N=28 (K dwarf 3:2).}
\end{figure}

\subsection{Stability-Filtered Null}

The stability-filtered null (Table~\ref{tab:null}; Figure~\ref{fig:null})
reveals that the G-dwarf bond fraction is indistinguishable from random packing
within stability bounds. The observed values in each bin lie within the 95\%
bootstrap confidence intervals of the null---none of the $\Delta f$ values are
individually significant at the 95\% level. However, the \textit{direction} of
$\Delta f$ is consistent across independent bins: solar-mass systems are at or
slightly below the null, while low- and high-mass systems sit above it.

The key result is not that any single bin shows significant excess. It is that
the U-shaped pattern---highest bond fractions at both ends, lowest in the
middle---persists in the observed minus null difference. If there were no
stellar-mass dependence, we would expect $\Delta f$ to scatter randomly around
zero. Instead, the three independent low+high+middle bins have $\Delta f$ of
$+2.5\%$, $+3.5\%$, and $-2.8\%$---a pattern consistent with a real U-shape
even if individual bins lack formal significance.

To test the U-shape explicitly, we perform a quadratic regression of
system-level bond fraction against stellar mass. The quadratic coefficient is
positive (concave-up, as expected for a U-shaped dependence) with $p = 0.06$,
suggesting the U-shape is not statistically significant at 95\% confidence but
is directionally consistent with the physical expectation. As an additional test, we permute host-star
masses across systems (10,000 iterations) while preserving the observed bond
fraction distribution; the observed U-shaped $\Delta f$ ordering---higher at
both ends than in the middle---occurs in only 11\% of permutations,
indicating that the pattern is unlikely to arise by chance under a
mass-independent null. We report these values as suggestive trends warranting
confirmation with larger samples, rather than as significant detections.

This means:
\begin{enumerate}
  \item G dwarfs are near the dynamical floor: their commensurability
        structure is consistent with what random survival alone would produce.
  \item The excess near-commensurability in M and F dwarfs, while not highly
        significant in individual bins, is directionally consistent across bins
        and strengthened by the U-shape pattern.
  \item Architectural excess---the apparent excess of near-commensurability
        beyond random survival---appears to be more pronounced around
        certain stellar types, though larger samples are needed to confirm.
\end{enumerate}

\begin{table}
\centering
\caption{Observed vs.\ stability-filtered null bond fractions across mass bins.
The null uses the same stellar mass bins as Table~\ref{tab:ushape} for consistency.
(We also verify results using broader bins---merging K and G dwarfs---to increase
null sample sizes within each bin; the pattern is unchanged.) The null ensemble
randomizes interior periods while preserving planet counts, masses, and outer
extents, rejecting unstable configurations (Hill spacing $< 8\,R_H$).
$\Delta f =$ observed minus null.}
\label{tab:null}
\begin{tabular}{lcccc}
\hline
Host bin & $N_{\rm sys}$ & Observed & Null & $\Delta f$\\
\hline
M dwarf ($<0.6$) & 47 & 66.0\% [56.6, 75.0] & 63.5\% & $+2.5\%$\\
K dwarf (0.6--0.9) & 125 & 54.6\% [48.5, 60.8] & 55.8\% & $-1.2\%$\\
G dwarf (0.9--1.1) & 113 & 47.0\% [40.4, 53.6] & 49.8\% & $-2.8\%$\\
F dwarf ($1.1$--$1.6$) & 47 & 56.0\% [45.3, 66.5] & 52.5\% & $+3.5\%$\\
\hline
\end{tabular}
\end{table}

\begin{figure}
\includegraphics[width=\columnwidth]{fig3_stability_null.png}
\caption{Observed bond fraction (system-level mean $\pm$ 95\% bootstrap CI)
compared to the stability-filtered null. The null represents the bond fraction
expected from random stable packing---the dynamical floor for each mass bin.
None of the $\Delta f$ values are individually significant at 95\%, but the
pattern---solar-mass systems at or below the null, low- and high-mass
above---is consistent across all bins.}
\label{fig:null}
\par\medskip\noindent{\small\itshape Alt text: Observed bond fractions with 95% bootstrap confidence intervals compared to the stability-filtered null (dynamical floor). None of the delta-f values (+2.5%, -1.2%, -2.8%, +3.5%) are individually significant at 95% confidence, but the U-shaped pattern (solar-mass at/below null, extremes above) is consistent across all bins.}
\end{figure}

\subsection{Notable Systems}

\textbf{TRAPPIST-1} (M dwarf, $0.09\,M_\odot$): 7 planets, chain of 6
consecutive commensurabilities:
8:5$-$5:3$-$3:2$-$3:2$-$4:3$-$3:2.

\textbf{HD~110067} (K dwarf, $0.80\,M_\odot$): 6 planets, chain of 5. A
near-uniform crystal: 3:2$-$3:2$-$3:2$-$4:3$-$4:3.

\textbf{TOI-1136} (G dwarf): 6 planets, mixed chain with 2:1, 7:5, and 3:2s.

\textbf{Kepler-223} (G dwarf): 4 planets, textbook
4:3$-$3:2$-$4:3 chain---a rare fully bonded G-dwarf system that illustrates
the exception within the broader trend.

\section{Two Candidate Mechanisms}
\label{sec:mechanisms}

The U-shape---higher bond fractions at both ends of the mass spectrum, lower
in the middle---demands an explanation. The pattern persists under
transit-only and higher-multiplicity filters, reducing the likelihood that
simple detection bias drives the result. Disk evolution models offer a coherent
heuristic framework. The chemical analogies in the subsections below are
descriptive tools, not literal claims about thermochemical processes. The
mechanisms are candidate interpretations rather than demonstrated processes.

\subsection{M Dwarf: Slow Cooking}

M-dwarf disks are often inferred to persist longer (5--10~Myr, compared to
2--5~Myr for G dwarfs)\citep{yasui2010,haisch2001}. Combined with low disk
masses, this produces slow planetary migration. Planets approach resonances
gradually, giving high capture probability. The extended disk lifetime also
provides post-capture dynamical settling: millions of years for orbits to
settle deeper into resonance.

Consistent with this picture, M dwarfs show the highest bond fraction, the
tightest bonds, and the widest bond palette (2:1, 5:3, 7:3, and 8:3 all
prominent).

\subsection{F Dwarf: Fast Chemistry}

F dwarfs have short-lived (1--3~Myr) but massive
disks\citep{sicilia_aguilar}. Migration is fast and vigorous. Planets sweep
through multiple resonances rapidly, giving capture opportunities through a
different channel: kinetic capture rather than slow settling. The short disk
lifetime leaves less time for post-capture settling, so bonds are looser on
average.

Consistent with this picture, F dwarfs have a bond fraction comparable to M
dwarfs (56.0\% vs.\ 66.0\%), but their bonds are consistently looser (2.09\%
deviation for 2:1 vs.\ M dwarfs' 1.76\%; 1.46\% vs.\ 1.07\% for 3:2).

\subsection{G Dwarf: The Dynamic Floor}

G dwarfs fall in the intermediate regime: disk lifetimes of 2--5~Myr, moderate
disk masses, moderate migration speed. Their bond fraction is consistent with
the stability-filtered null, placing them at or near the dynamical floor. The
U-shape may be a general feature of systems where bond formation requires a
specific dynamical regime---either slow enough for adiabatic capture (M dwarfs)
or fast enough for multiple kinetic capture events (F dwarfs)---and G dwarfs
fall in the intermediate range where neither mechanism operates efficiently.

\subsection{Age Correlations}
\label{sec:age_correlations}

\begin{figure}
\includegraphics[width=\columnwidth]{fig4_age_bond.png}
\caption{System-level bond fraction as a function of isochronal age, by stellar
type. The U-shape in bond fraction across mass is mirrored by opposite age
correlations: M dwarfs show tentatively increasing bond fraction with age
($+0.8\%$ Gyr$^{-1}$, or approximately $+8\%$ over 10 Gyr); K, G, and F dwarfs all show decreasing bond fraction with age.}
\label{fig:agebond}
\par\medskip\noindent{\small\itshape Alt text: Four-panel scatter plot of system-level bond fraction vs isochronal age, by stellar type. M dwarf (N=29): slope +0.8% Gyr^{-1}, p=0.39. K dwarf (N=63): slope -1.5% Gyr^{-1}, p=0.03. G dwarf (N=78): slope -0.6% Gyr^{-1}, p=0.16. F dwarf (N=21): slope -2.0% Gyr^{-1}, p=0.26. Only the K-dwarf trend is individually significant (p=0.03) before Bonferroni correction (corrected threshold p<0.0125).}
\end{figure}

The results (Table~\ref{tab:agebond}, Figure~\ref{fig:agebond}) show a
striking pattern, though the uncertainties are large. Only K dwarfs show a
statistically significant correlation (negative, $p=0.03$). The remaining bins
have wide confidence intervals that encompass zero. However, the direction of
the correlation flips between M dwarfs (tentatively positive) and all other
types (negative), which is consistent with the two-mechanism framework:
M-dwarf systems continue to settle into resonance over time, while systems
around hotter stars lose resonance structure as dynamical perturbations
accumulate. We note that with four age correlation tests, a single $p=0.03$ is
not significant after a Bonferroni correction (adjusted threshold $p<0.0125$);
we report it as suggestive given the directional consistency across bins.

\begin{table}
\centering
\caption{Linear fit of bond fraction vs.\ isochronal age, by stellar type.}
\label{tab:agebond}
\begin{tabular}{lccc}
\hline
Star type & $N_{\rm sys}$ & Slope (\% Gyr$^{-1}$) & $p$-value\\
\hline
M dwarf & 29 & $+0.8 \pm 0.9$ & 0.39\\
K dwarf & 63 & $-1.5 \pm 0.7$ & 0.03\\
G dwarf & 78 & $-0.6 \pm 0.4$ & 0.16\\
F dwarf & 21 & $-2.0 \pm 1.7$ & 0.26\\
\hline
\end{tabular}
\end{table}

\section{Predictions and Tests}

The two-mechanism framework generates specific, testable predictions distinct
from a null hypothesis of uniform bond formation:
\begin{itemize}
  \item The M-dwarf age-bond correlation ($+0.8\%$ Gyr$^{-1}$; approximately
        $+8\%$ over a Hubble time,
        Section~\ref{sec:age_correlations}) should strengthen with larger, more
        homogeneous age samples: older M dwarfs show tighter bonds, consistent
        with continued dynamical settling.
  \item High-resolution ALMA observations of M-dwarf disks should confirm
        longer dissipation timescales relative to FGK stars.
  \item Bond fraction around F dwarfs should anticorrelate with disk accretion
        rate (faster accretion $\rightarrow$ less settling).
  \item Future high-completeness surveys of M-dwarf multi-planet systems should
        preferentially recover compact near-commensurable chains.
\end{itemize}

\section{Solar System within the Framework}

The Solar System provides a single, well-characterized data point at the
G-dwarf location---the bottom of the U-shape. Among adjacent planets from
Mercury to Neptune, only one pair (Venus--Earth, $P_\oplus/P_\venus \approx
365.25/224.7 \approx 1.626$, within 0.06\% of 13:8) falls
within the $\delta \leq 0.03$ bond criterion. Pluto-Neptune (3:2 at $\sim
350$~AU) is a clean commensurability but involves a dwarf planet at extreme
orbital separation. The Solar System is one of the least bonded planetary
architectures known, consistent with its G dwarf host at the minimum of the
bonding curve.

\section{Acknowledgements}

The author used AI language models (Claude, DeepSeek, and ChatGPT) for
drafting assistance, methodological discussion, and code-review support.
The author takes full responsibility for the analysis, interpretation, and
final text. We thank an anonymous colleague for a rigorous reading of an early
draft and for suggestions that strengthened the methodology and presentation.
The author has no funding sources to declare. This research has made use of
the NASA Exoplanet Archive, which is operated by the California Institute of
Technology, under contract with the National Aeronautics and Space
Administration under the Exoplanet Exploration Program.

\section*{Data Availability}

All data used in this study are publicly available from the NASA Exoplanet
Archive at \url{https://exoplanetarchive.ipac.caltech.edu}. The
stability-filtered null simulation code is available from the corresponding
author upon reasonable request.

\bibliographystyle{mnras}
\bibliography{references}

\label{lastpage}

\end{document}
